\section{行为主体 - 单位犯罪}
\begin{enumerate}
  \item 分类：
        \begin{enumerate}
          \item 纯正的：只能由单位构成。
          \item 不纯正：单位 + 个人。生产、销售伪劣产品。
          \item 纯正的自然人：只能由自然人构成。诈骗罪、盗窃罪、金融犯罪：
                \begin{enumerate}
                  \item \zd{贷款}诈骗、\zd{信用}卡诈骗、\zd{有价}证券诈骗（没考过）\yjt 只能自然人。
                  \item 集资、信用证、保险诈骗 \yjt 单位和自然人。
                \end{enumerate}
        \end{enumerate}
  \item 主体条件。
        \begin{enumerate}
          \item 一般情况下，不要求单位有法人资格，\zd{私营企业}需要。
          \item 合伙不能构成单位犯罪（无法人资格），按自然人处理。
          \item 子公司可以，具有法人资格。
          \item 分公司、分支机构、内设机构：\zd{以自己名义去犯罪 + 违法所得归机构所有}。
        \end{enumerate}
  \item 主观条件。\zd{体现单位的意志}。
        \begin{enumerate}
          \item 集体/依职权决策。
          \item 既可以是故意犯罪，也可以是过失（工程重大安全事故）。
          \item \qd{以犯罪目的成立} \hz \zd{中途变更目的}为犯罪 \yjt 按自然人犯罪处理。
        \end{enumerate}
  \item 单位犯罪和个人犯罪的关系。
        \begin{enumerate}
          \item \bt{单位犯罪和单位内部人员犯罪的区分}。
                \begin{enumerate}
                  \item 主观：单位意志还是个人意志。
                  \item 客观：\zd{为谁}牟取利益。
                        \begin{enumerate}
                          \item 主观单位，客观个人（\zd{少数特定}人不能算为集体）\yjt 个人犯罪。
                          \item 主观个人，客观单位 \yjt 个人。\al{殴打执行法官}，个人意志，客观单位利益，但\zd{个人行为无授权}。
                        \end{enumerate}
                \end{enumerate}
          \item \zd{不是各成员间的共同犯罪}（体现的是个人的意志）。
          \item \zd{单位实施纯正的自然人犯罪}：\al{集体决定盗电}。\zd{定单位直接责任人}。
        \end{enumerate}
  \item 处罚。
        \begin{enumerate}
          \item 原则双罚制，既罚单位，也罚个人。例外：单罚，只罚个人，不罚单位。
          \item 对单位的处罚\zd{只有罚金}；个人包括主管人员和直接责任人。
          \item 单位真没了 \yjt 直接追究直接责任人的刑事责任；单位假没了 \yjt 追究\zd{原单位}。
        \end{enumerate}
\end{enumerate}
